\subsection{拟不变测度}

命题 11. --- 设 G 为局部紧群，\(\mu\) 为 G 上的左哈尔测度。G 上的测度 \(\nu \neq 0\) 左拟不变，当且仅当 \(\nu\) 等价于 \(\mu\)。

充分性显然。设 \(\nu \neq 0\) 是左拟不变测度，我们证明 \(\nu\) 等价于 \(\mu\)。可以限于 \(\nu > 0\) 的情形。令 A 为 G 的紧子集。我们将证明（这也就证明了命题），条件 \(\mu(A) = 0\) 与 \(\nu(A) = 0\) 等价（第五章，§5，第 5 号，定理 2）。

\begin{enumerate}
\def\labelenumi{\alph{enumi})}
\tightlist
\item
  对于每个 \(f \in \mathcal{K}_{+}(\mathrm{G})\)，函数 \((x, y) \mapsto f(x)\varphi_{\mathrm{A}}(xy)\) 在 \(G \times G\) 上是 \((\nu \otimes \mu)\)-可积的，因为它是上半连续且有界的，并且其支集包含于紧集 \(K \times K^{-1}A\)，其中取 \(K = Supp f\)。因此，根据勒贝格--富比尼定理，
\end{enumerate}

\[
\int d \nu (y) \int \varphi_ {\mathrm{A}} (x y) f (x) d \mu (x) = \int f (x) d \mu (x) \int \varphi_ {\mathrm{A}} (x y) d \nu (y).\tag{36}
\]

\begin{enumerate}
\def\labelenumi{\alph{enumi})}
\setcounter{enumi}{1}
\tightlist
\item
  假设 \(\nu(\mathbf{A}) = 0\)。根据假设，\(\nu(x\mathbf{A}) = 0\) 对所有 \(x \in G\) 成立，故（36）的右端为零。因此存在 \(\nu\)-可忽略集 \(N_{f}\)，使得对于 \(y \notin N_{f}\)，
\end{enumerate}

\[
0 = \int \varphi_ {\mathrm{A}} (x y) f (x) d \mu (x) = \Delta_ {\mathrm{G}} (y) ^ {- 1} \int \varphi_ {\mathrm{A}} (x) f (x y ^ {- 1}) d \mu (x).\tag{37}
\]

令 B 为 G 的紧子集，使 \(\nu(\mathrm{B})\neq0\)，并取 \(\mathcal{K}_{+}(\mathrm{G})\) 中在 \(AB^{-1}\) 上等于 1 的函数 f。于是存在 \(y\in B\) 使（37）成立。但由于 \(\varphi_{\mathrm{A}}(x)f(xy^{-1})=\varphi_{\mathrm{A}}(x)\) 对于 \(y\in B\) 成立，这证明 \(\mu(\mathrm{A})=0\)。

\begin{enumerate}
\def\labelenumi{\alph{enumi})}
\setcounter{enumi}{2}
\tightlist
\item
  假设 \(\mu(A) = 0\)。那么，对于每个 \(f \in \mathcal{K}_{+}(G)\)，（36）的左端为零，从而右端也为零。因此存在局部 \(\mu\)-可忽略集 \(M\)，使 \(\int \varphi_{A}(xy) d\nu(y) = 0\) 对 \(x \notin M\) 成立。由于 \(\mu \neq 0\)，可知 \(\nu(xA) = 0\) 对某个 \(x \in G\) 成立，从而 \(\nu(A) = 0\)。
\end{enumerate}

将命题 11 应用于 \(G^{0}\)，可见右拟不变测度与左拟不变测度相同。它们统称 G 上的拟不变测度。

